%!TEX root = ../../thesis-trier.tex

In \cref{prelim:autsyncw:prop:two-sets-min-rank}
we prove, using \cref{sync_max:prop:two-sets-containing-2set},
that distinguishability of pairs of $k$-sets in automata with at least $k+1$ states 
implies synchronizability for $k \ge 3$, and for $k = 2$
existence of a word of rank at most two.
Then we list several consequences, for example,
that a strongly connected automaton in which pairs of $2$-sets
are distinguishable is synchronizing.
\Cref{sync_max:ex:2-sets-binary,sync_max:ex:2-sets}
give non-synchronizing automata where all pairs of $2$-sets
are distinguishable. Lastly, in~\cref{sync_max:prop:disjointklsets}, 
for automata whose transformation monoid contains a transitive group,
we give a sufficient condition for synchronizability in terms 
of the distinguishability of certain disjoint sets.

\begin{lemma}
	\label{sync_max:prop:two-sets-containing-2set}
	Let $\semiautomat{A}{Q}{\Sigma}{\delta}$.
	Suppose for all distinct states $p, q \in Q$
	there exist distinguishable subsets $A, B \subseteq Q$
	with $\{p,q\} \subseteq A \cap B$.
	Then $\mathcal A$ is synchronizing.	
\end{lemma}
\begin{proof}
	A word $u \in \Sigma^*$ that distinguishes $A$ from $B$
	maps at least one subset to a singleton subset,
	which implies $\applytrans{\delta}{p}{u} = \applytrans{\delta}{q}{u}$.
	As this holds for any pair of distinct states $p,q \in Q$, by \cref{prelim:thm:pair_sync_condition} 
	the automaton $\mathcal A$ is synchronizing.    			
\end{proof}

\begin{proposition} \label{prelim:autsyncw:prop:two-sets-min-rank} 	\label{sync_max:cor:k-ge-3-sync}
	Let $\semiautomat{A}{Q}{\Sigma}{\delta}$ and $2 \le k,l \le |Q|$ with $|Q| \ge k + 1$
	or $|Q| \ge l + 1$.
	Suppose each $k$-set is distinguishable from
	each distinct $l$-set of states in~$\mathcal A$
	and let $u$ be a word of minimal rank~$r$.
	For $k \ge 3$ or $l \ge 3$, we have $r = 1$, i.e., $\mathcal A$ is synchronizing.
	For $k = l = 2$, we have~$r \le 2$ where for $r = 2$
	the set $\delta(Q, u)$ is permuted by each letter, and so
	a union of sink components by \cref{prelim:autsyncw:lem:perm-each-letter-sink}.
\end{proposition}
\begin{proof}	
	Without loss of generality, let $k \le l$.
	First, suppose $l \ge 3$.
	Let $p,q \in Q$ be distinct.
	By assumption, there exist distinct $p_3, \ldots, p_k, p_{k+1} \in Q \setminus \{p,q\}$.
	Set $A = \{p,q,p_3, \ldots, p_k\}$.
	Next, there also exists an $l$-subset $B \subseteq Q$
	with $\{p,q,p_{k+1}\} \subseteq B$.
	Since $p_{k+1} \notin A$
	we have $A \ne B$.
	By assumption, $A$ and~$B$ are distinguishable.
	The, from \cref{sync_max:prop:two-sets-containing-2set} we conclude that $r = 1$. 
				
	Next, suppose $k = l = 2$.
	Assume $r \ge 3$. Then we can pick three distinct states $p,q,s \in \delta(Q, u)$.
	By assumption, the sets $\{p,q\}$ and $\{p,s\}$
	are distinguishable by some word $v \in \Sigma^*$.
	Since one of both sets is mapped to a singleton,
	we find $|\delta(Q, uv)| < |\delta(Q, u)|$, contradicting
	the choice of~$u$ as a word of minimal rank.
	So, we have $r \le 2$.
	Now, suppose $r = 2$.
	Assume there exists $a \in \Sigma$
	such that $\delta(\delta(Q, u), a) \ne \delta(Q, u)$.
	Since~$r$ is the minimal possible rank, $\delta(\delta(Q, u), a)$
	is a $2$-set.
	Then, by assumption, there exists a word $v$
	distinguishing $\delta(\delta(Q, u), a)$
	and $\delta(Q, u)$.
	However, then one of both subsets is mapped
	to a singleton, which is not possible
	since the minimal possible rank of a word is $2$.
	Hence, $\delta(\delta(Q, u), a) = \delta(Q, u)$
	for each $a \in \Sigma$.	
\end{proof}

Next, we list a few
consequences of \cref{sync_max:prop:two-sets-containing-2set,sync_max:cor:k-ge-3-sync}.

\begin{corollary} \label{sync_max:cor:two-sets-containing-2set}
	Let $\semiautomat{A}{Q}{\Sigma}{\delta}$ be an $n$-state automaton and $2 \le k < n$.
	
	\begin{enumerate}
		\item \label{sync_max:cor:two-sets-containing-2set:1} If $n \ge 4$ and all pairs of distinct $3$-sets
		are distinguishable in $\mathcal A$, then $\mathcal A$ is synchronizing.
		
		
		\item \label{sync_max:cor:two-sets-containing-2set:2} If for every $k$-set $A \subseteq Q$ there exists $B \subseteq Q$
		with $A \subseteq B$ such that $A$ and~$B$ are distinguishable,
		then $\mathcal A$ is synchronizing.
		
		\item \label{sync_max:cor:two-sets-containing-2set:item:kl}  Let $0 < l \le n - k$.
		If every $k$-set is distinguishable from 
		every $(k+l)$-set, then $\mathcal A$ is synchronizing.  
		
		\item \label{sync_max:cor:two-sets-containing-2set:4} If every $k$-set is distinguishable from the whole
		state set, then the automaton is synchronizing.
	\end{enumerate}	
\end{corollary}
\begin{proof}
	\ref{sync_max:cor:two-sets-containing-2set:1}: This follows from~\cref{sync_max:cor:k-ge-3-sync} with $k = 3$.
	
	\ref{sync_max:cor:two-sets-containing-2set:2}: For distinct $p,q \in Q$ choose some $k$-set with $p, q \in A$.
	By assumption, there exists $B \subseteq Q$ containing $A$ 
	so that $A$ and $B$ are distinguishable.
	Then \cref{sync_max:prop:two-sets-containing-2set} gives the claim.
	
	\ref{sync_max:cor:two-sets-containing-2set:item:kl} and~\ref{sync_max:cor:two-sets-containing-2set:4} clearly follow from~\ref{sync_max:cor:two-sets-containing-2set:2}.
\end{proof}

\begin{corollary} \label{prelim:aut-sync-words:subset:cor:no-letter-fixes-two-sets-and-dist-2-sets-implies-sync} \label{sync_max:prop:2-sets-dist-SCC-implies-sync}
Suppose $\semiautomat{A}{Q}{\Sigma}{\delta}$ is an automaton with at least three states in which all pairs of distinct $2$-sets
are distinguishable.
\begin{enumerate}
\item \label{sync_max:prop:2-sets-dist-SCC-implies-sync:1} If for distinct $p,q \in Q$ 
we have $\delta(\{p,q\}, a) \ne \{p,q\}$ for some $a \in \Sigma$,
then $\mathcal A$ is synchronizing.

\item \label{sync_max:prop:2-sets-dist-SCC-implies-sync:2} If~$\mathcal A$ is strongly connected, then it is synchronizing.

\item \label{sync_max:prop:2-sets-dist-SCC-implies-sync:3} Either $\mathcal A$ has a unique sink component and is synchronizing, or it has 
precisely two sink components each of size one.
\end{enumerate}
\end{corollary}
\begin{proof}
\ref{sync_max:prop:2-sets-dist-SCC-implies-sync:1}: Let $u \in \Sigma^*$ be a word of minimal possible rank.
By \cref{prelim:autsyncw:prop:two-sets-min-rank} 
we have $|\delta(Q, u)| \le 2$
where $|\delta(Q, u)| = 2$ is not possible
since otherwise every letter would permute this set
by \cref{prelim:autsyncw:prop:two-sets-min-rank}.

\ref{sync_max:prop:2-sets-dist-SCC-implies-sync:2}: Let~$u$ be a word of minimal rank
and assume $|\delta(Q, u)| \ge 2$.
By \cref{prelim:autsyncw:prop:two-sets-min-rank},
$\delta(Q, u)$ has size two and equals a union of sink components.
Since~$\mathcal A$ is strongly connected, $Q$ is its only sink component
and  we find $\delta(Q, u) = Q$,
which is not possible since~$|Q| \ge 3$.

\ref{sync_max:prop:2-sets-dist-SCC-implies-sync:3}:	By \cref{prelim:autsyncw:lem:minrank}, $\mathcal A$
has at most two sink components.
First, if it has a unique sink component~$S \subseteq Q$,
then $\mathcal A_{|S}$ is strongly connected
and all pairs of distinct $2$-sets from~$S$ are distinguishable.
Thus, by Item~\ref{sync_max:prop:2-sets-dist-SCC-implies-sync:2}, $\mathcal A_{|S}$ is synchronizing.
Since the set of maximal states equals the union of all sink components
by \cref{prelim:order:maximal-elements-upper-set}, $S$ equals the set of maximal states
and so~$\mathcal A$ is synchronizing by \cref{prelim:thm:SCC_sync_words}.
Next, suppose~$\mathcal A$ has two sink components.
Then it is not synchronizing by \cref{prelim:autsyncw:lem:minrank}.
By \cref{prelim:autsyncw:prop:two-sets-min-rank},
$\delta(Q, u)$ for a word of minimal rank~$u$
equals a union of sink components and has size two.
Since this set has non-empty intersection with both sink components by \cref{sync_max:eqn:upper_set}, 
both sink components must have size one.
\end{proof}

Next, \cref{sync_max:ex:2-sets-binary,sync_max:ex:2-sets} give infinite families of non-synchronizing automata 
where all pairs of $2$-sets are distinguishable, demonstrating that the case $k = 2$ in
\cref{prelim:autsyncw:prop:two-sets-min-rank} is special in the sense that
it does not imply synchronizability similar to the case $k > 2$.
By adding in these examples another letter that swaps the two maximal states and fixes every other state, we obtain non-synchronizing automata 
with a unique sink component where all pairs of $2$-sets are distinguishable.

\begin{example} 
	\label{sync_max:ex:2-sets}
	Set $\semiautomat{A_n}{[n]}{\Sigma}{\delta}$
	for $n \ge 3$
	where $\Sigma = \{ a_{ij} \mid i \in \{0,1\}, j \in \{2,\ldots,n-1\}\}$
	and, for $i \in \{0,1\}$, $j \in \{2,\ldots,n-1\}$,
	we have
	\[
	\delta(k, a_{ij}) = \begin{cases}
		i & \text{if } k = j, \\
		k & \text{otherwise,}
	\end{cases} 
	\]
	for $k \in [n]$.
	See \cref{sync_max:fig:2-sets}
	for~$\mathcal A_3$.
	Now, consider two distinct $2$-sets $\{p, q\}$
	and $\{s, t\}$ of states. First, without loss of generality,
	suppose $\{p, q\} \subseteq \{2,\ldots,n-1\}$.
	Then $\delta(p, u) = \delta(q, u) = 0$
	and $|\delta(\{s,t\}, u)| = 2$
	for $u = a_{0p} a_{0q}$.
	Otherwise, $\{0,1\}$ has non-empty intersection
	with both $2$-sets where we assume, without loss of 
	generality, $p = 0$ and $q \ne 1$.
	Then,
	$\delta(p, a_{0q}) = \delta(q, a_{0q}) = 0$
	and $|\delta(\{s, t\}, a_{0q})| = 2$.
\end{example}

\begin{example}
	\label{sync_max:ex:2-sets-binary}
	Let $\semiautomat{A_n}{Q_n}{\{a,b\}}{\delta_n}$ be the automaton given by the state diagram in \cref{sync_max:fig:2-sets-binary} where $n \ge 3$. Consider two distinct sets of states  $\{p_1, p_2\}, \{q_1, q_2\}$ with $p_1 < p_2$ and $q_1 < q_2$. %Note that this implies $p_1, q_1 \le n - 2$.
	First, without loss of generality, suppose $p_1 < q_1$.
	In this case, for $q_2 \le n - 2$, we have $\delta(\{q_1, q_2\}, a^{n - 2 - q_1}) = \{n - 2\}$ 
	and $|\delta(\{p_1, p_2\}, a^{n - 2 - q_1})| = 2$.
	For $q_2 = n - 1$ and $q_1 \ne n - 2$,
	we have $\delta(\{q_1, q_2\}, b^{n - 2 - q_1}) = \{n-1\}$
	and $|\delta(\{p_1, p_2\}, b^{n-2 - q_1})| = 2$.
	For $q_2 = n - 1$ and $q_1 = n - 2$,
	if $p_2 = n - 1$ the word $b^n$ distinguishes both subsets,
	and if $p_2 \le n - 2$ the word $a^n$ distinguishes both subsets.
	Next, suppose $p_1 = q_1$. In this case, without loss of generality,
	assume $p_2 < q_2$. % p_2 auf n - 2
	For $q_2 = n - 1$, we have $\delta(\{p_1, p_2\}, a^n) = \{n-2\}$
	and $\delta(\{q_1, q_2\}, a^n) = \{n - 1, n - 2\}$.
	Otherwise, for $q_2 < n - 1$,
	we have $\delta(\{p_1, p_2\}, u) = \{ n - 2 \}$
	and $\delta(\{q_1, q_2\}, u) = \{ n - 1, n - 2\}$
	for $u = a^{n - 3 - q_2}ba^n$. 
\end{example}

\todo[inline]{maßangaben auf em umstellen}
\begin{figure}
	\centering
	\begin{subfigure}[b]{0.44\textwidth}		
		\centering
		\begin{tikzpicture}[node distance = 2.5cm and 1.25cm]
			
			\node[state] (3) {$3$};
			
			\node[state, left = of 3] (4) {$4$};
			\node[state, right = of 3] (2) {$2$};
			\node[state, above = of 4] (0) {$0$};
			\node[state, above = of 2, yshift=-0.0cm] (1) {$1$};
			
			\path[->] (2) edge [bend left=10,above,sloped,pos=.7] node {$a_{02}$} (0)
			(2) edge [right] node {$a_{12}$} (1); 
			
			\path[->] (3) edge [bend left=15, above, sloped, pos=0.49] node {$a_{03}$} (0)
			(3) edge [bend right=15, right] node {$a_{13}$} (1)
			(4) edge [above, sloped] node {$a_{04}$} (0)
			(4) edge [bend right=15, pos=0.7,above, inner sep=9pt] node {$a_{14}$} (1);	
		\end{tikzpicture}	
		\caption[Non-Synchronizing and pairs of $2$-sets distinguishable]{Automaton as in \cref{sync_max:ex:2-sets}.}
		\label{sync_max:fig:2-sets}
	\end{subfigure}
	\hfill
	\begin{subfigure}[b]{0.55\textwidth}
		\centering
		\begin{tikzpicture}[inner sep=0pt, minimum size=20pt, node distance = 0.6cm and 0.7cm]
			\node[state] (0) {$0$};
			\node[state, right = of 0] (1) {$1$};
			\node[right = of 1] (dots) {$\cdots$};
			\node[state, right = of dots] (n-3) {$n-3$};
			\node[state, right = of n-3] (n-2) {$n-2$};
			\node[state, below = of n-2] (n-1) {$n-1$};
			
			\path[->] (0) edge [above] node {$a,b$} (1)
			(1) edge [above] node {$a,b$} (dots)
			(dots) edge [above] node {$a,b$} (n-3)
			(n-3) edge [above] node {$a$} (n-2)
			(n-3) edge [below] node {$b$} (n-1);
		\end{tikzpicture}
		\caption[Binary non-Synchronizing and pairs of $2$-sets distinguishable]{Binary automaton as in \cref{sync_max:ex:2-sets-binary}.} 
		\label{sync_max:fig:2-sets-binary}
	\end{subfigure}
	\caption[Tow non-synchronizing automata where all pairs of $2$-sets distinguishable]{Two non-synchronizing automata where all pairs of distinct $2$-sets are distinguishable, self-loops omitted.}
	\label{sync_max:fig:non-sync-two-sets-dist}
\end{figure}

Lastly, we consider the case that the transformation monoid contains a transitive group. In this case, we state a sufficient condition for synchronizability
in terms of the distinguishability of certain disjoint subsets. We need this result later in a more algebraic context.
For this reason, we give a formulation in terms of transformation monoids without reference to automata.

\begin{proposition} 
	\label{sync_max:prop:disjointklsets}
	Let $M \le_M \mathcal T_n$ be a transformation
	monoid with a transitive group of units
	and $2 \le k, l < n$ such that $kl < n$.
	If every $k$-set is distinguishable
	from every disjoint $l$-set,
	then $M$ contains a constant map. 		
\end{proposition}
\begin{proof}
	Without loss of generality, assume $k \le l$.
	Let $p, q \in [n]$ be two distinct states.
	Pick $A,B \subseteq [n]$ with $k = |A|$, $l = |B|$
	and $p,q \in A \cap B$ (note that the following argument includes the case $k = l = 2$ and $A = B = \{p,q\}$).
	As $G$ is transitive, by \cref{thm:separating_neumann} 
	we have $\applyfunc{A}{g} \cap B = \emptyset$
	for some $g \in G$.
	By assumption there exists $h \in M$
	mapping precisely one of both subsets $\applyfunc{A}{g}$
	or $B$ to a singleton subset.
	If $|\applyfunc{\applyfunc{A}{g}}{h}| = 1$,
	then, as $\{p,q\} \subseteq A$,
	the transformation $\concat{g}{h}$
	maps $\{p,q\}$ to a singleton.
	If $|\applyfunc{B}{h}| = 1$,
	then, as $\{p,q\} \subseteq B$,
	the transformation~$h$ maps $\{p,q\}$
	to a singleton.	
	Using the correspondence between automata
	and transformation semigroups (see \cref{prelim:remarks_trans_semigrp_semiautomata}),
	the pair criterion (\cref{prelim:thm:pair_sync_condition})
	implies that $M$ contains a constant map as all pairs of states $p,q$ 
	can be mapped to the same state.	
\end{proof}